% Section 1: Abstract
% Target length: 150-250 words (ArXiv convention: keep under 250)
% Status: DRAFT - mirrors abstract.txt, needs tightening once Sections 4-7 are final.

\begin{abstract}
Automated textbook generation with large language models typically allocates
a flat word budget per chapter or per concept, independent of how important
each concept is to the rest of the curriculum. We report an informal but
consistent observation across two AI-assisted textbook projects: generated
chapter quality improved, by subjective reading, when the word budget for
each concept was scaled by the concept's in-degree in the course's learning
graph --- the number of other concepts that directly depend on it. We show
that raw in-degree undercounts concepts that are foundational only
\emph{transitively} (few direct dependents but many indirect ones), and
propose the Concept Impact Score (CIS), a PageRank-style recursive importance
measure over the concept dependency graph. Because a learning graph is
acyclic by construction, CIS reduces to an exact, closed-form, single-pass
computation over a topological ordering, requiring no damping factor. We
formalize the chapter-generation problem as three separable sub-problems ---
concept-to-chapter partitioning, within-chapter concept ordering, and
per-concept content-budget allocation --- and propose CIS-weighted
allocation for the last of these. On a real 200-concept, 277-edge Algebra~I
learning graph, CIS and in-degree rank concepts substantially differently;
several low-in-degree concepts (e.g. \emph{Constant}, \emph{Coefficient})
rank in the top five by CIS due to transitive reach. We describe a planned
blind-comparison study across uniform, in-degree-weighted, and CIS-weighted
conditions at fixed total chapter length to test the hypothesis directly.

% TODO: finalize word count to <=250 words after Sections 4-7 are drafted.
\end{abstract}
